\subsection{Matriz Indicadora de Módulos}
\label{sec:mim}

A MIM foi elaborada para a C1 antes da prioridade dada à fronteira M1. A matriz reúne 12 direcionadores e 13 portadores, com OP-01 fora da base. Os 156 escores e suas somas são preservados como histórico de julgamentos de projeto, sem selecionar C1. A descrição dos portadores foi adaptada para a peça descoberta, mas isso não revalida os escores. A mudança exige reavaliar os motivos de modularização, as interfaces e os agrupamentos antes de adotá-los para M1. A planilha sinaliza essa pendência em todas as relações do registro.

A MIM indica motivos para modularizar e candidatos a agrupamento; interfaces e acoplamentos também precisam ser examinados \citep{erixon1998,blees2009}. A escala adotada é 0 (sem motivo identificado), 1 (fraco), 3 (médio), 5 (forte), conforme a legenda do material-base da disciplina \citep{basemim}. A literatura consultada apresenta forte como 9; a convenção da base foi mantida explicitamente, sem atribuir-lhe a escala clássica.

\begin{table}[H]
\centering\scriptsize\caption{Direcionadores e critérios de julgamento da MIM}\label{tab:mim-drivers}
\begin{tabularx}{\textwidth}{L{0.8cm}L{4.5cm}Y}
\toprule\textbf{ID}&\textbf{Direcionador}&\textbf{Critério}\\\midrule
D01 & Carry-over & Reaproveitar o projeto em futuras gerações; não é reuso por ciclos. \\
D02 & Evolução tecnológica & Isolar partes sujeitas a mudança de tecnologia. \\
D03 & Alterações planejadas de projeto & Separar mudanças previstas no plano de produto. \\
D04 & Especificação técnica & Acomodar diferenças de requisitos entre variantes. \\
D05 & Estilo & Separar alterações de aparência ou identidade de marca. \\
D06 & Unidade comum & Usar a mesma parte física nas variantes atuais. \\
D07 & Processo e organização & Separar processos especializados ou unidades de trabalho. \\
D08 & Testes em separado & Caracterizar o portador antes do ensaio integrado. \\
D09 & Compra de produto pronto & Adquirir solução funcional pronta; matéria-prima não basta. \\
D10 & Manutenção e mantenabilidade & Acessar, limpar e repor partes durante a vida. \\
D11 & Atualização & Melhorar o produto em serviço, além de reparar desgaste. \\
D12 & Reciclagem & Separar materiais e facilitar a destinação. \\
\bottomrule\end{tabularx}
\fonte{Critérios de modularização, sintetizados de \citet{erixon1998}, \citet{blees2009} e da base; aplicação própria.}
\end{table}

Carry-over trata de reaproveitar projeto entre gerações, não de usar a mesma unidade por vinte ciclos. Compra de produto pronto não equivale a comprar polímero ou painel transparente. D03 ficou zerado porque não há roteiro de alterações futuras; D05, porque estilo não integra o recorte. Zero explicita ausência de motivo de modularização e não dispensa função ou requisito.

\begin{table}[H]
\centering\scriptsize\caption{Portadores propostos para C1 sob M1, com julgamentos a reavaliar}\label{tab:mim-portadores}
\begin{tabularx}{\textwidth}{L{0.8cm}L{3.7cm}YL{2.5cm}}
\toprule\textbf{ID}&\textbf{Função}&\textbf{Portador proposto}&\textbf{RP principais}\\\midrule
S01 & Amortecer impacto & Calços de EPP com assento e retenção a definir & RP1, RP12, RP13, RP18 \\
S02 & Imobilizar louça & Cintas com tensão limitada e apoios protegidos na cerâmica & RP2 \\
S03 & Acomodar variação & Guias deslizantes com curso e trava definidos & RP7 \\
S04 & Expor condição & Visor na barreira do próprio dispositivo, sobre peça descoberta & RP4, RP6 \\
S05 & Prover pega & Cavas de pega no corpo estrutural & RP3, RP5 \\
S06 & Portar louça & Interface estrutural para carrinho de obra & RP3, RP8, RP14 \\
S07 & Reter fragmentos & Bandeja e barreira envolvente contínua & RP6 \\
S08 & Suportar carga & Colunas e caminhos de carga do corpo polimérico & RP11, RP13 \\
S09 & Unir e liberar partes & Fechos sobre-centro com retenção secundária & RP17, RP20, RP22 \\
S10 & Comunicar uso e exibir identidade & Marcação de uso e código gravado na estrutura permanente & RP10, RP16, RP21 \\
S11 & Estabilizar conjunto & Geometria de base alargada sem rodízios livres & RP9 \\
S12 & Recolher partes no retorno & Laterais recolhíveis mantendo a base comum & RP14, RP17 \\
S13* & Registrar evento (OP-01) & Estudo de indicador comercial ou registrador & OP-01 \\
\bottomrule\end{tabularx}
\fonte{Função Global, morfológica e C1. S13* é opção de estudo, sem inclusão na referência.}
\end{table}

Manutenção, vida, reciclagem, troca e custo são critérios transversais (RP15, RP18, RP19, RP20 e RP23), examinados sobre os portadores. O catálogo distingue 17 funções: F15 viabiliza acesso e F19 separação material; esses suportes são distribuídos nas interfaces. S10 reúne uso e identidade (F16 e F21). Uniformidade qualifica F01 e umidade qualifica F11. S12 recolhe partes no vazio, sem comprovar RP14 no estado carregado. Sob M1, S04 deve permitir acesso real através da própria barreira, com a peça descoberta e sem depender de S13. Apoios e cintas exigem avaliação de pressão sobre a cerâmica. Os escores herdados não demonstram que essas novas interfaces estão resolvidas.

\begin{table}[H]
\centering\scriptsize\caption{Escores históricos da MIM da C1, sujeitos a reavaliação sob M1}\label{tab:mim-pontos}
\begingroup\setlength{\tabcolsep}{2pt}\renewcommand{\arraystretch}{1.25}
\begin{tabularx}{\textwidth}{L{3.8cm}*{13}{>{\centering\arraybackslash}X}}
\toprule\textbf{Direcionador}&\textbf{S01}&\textbf{S02}&\textbf{S03}&\textbf{S04}&\textbf{S05}&\textbf{S06}&\textbf{S07}&\textbf{S08}&\textbf{S09}&\textbf{S10}&\textbf{S11}&\textbf{S12}&\textbf{S13*}\\\midrule
Carry-over & 1 & 3 & 3 & 1 & 3 & 3 & 1 & 3 & 3 & 3 & 3 & 3 & 0 \\
Evolução tecnológica & 1 & 0 & 1 & 3 & 0 & 0 & 0 & 0 & 0 & 0 & 0 & 0 & 5 \\
Alterações planejadas de projeto & 0 & 0 & 0 & 0 & 0 & 0 & 0 & 0 & 0 & 0 & 0 & 0 & 0 \\
Especificação técnica & 5 & 3 & 5 & 3 & 1 & 1 & 1 & 1 & 1 & 0 & 1 & 3 & 3 \\
Estilo & 0 & 0 & 0 & 0 & 0 & 0 & 0 & 0 & 0 & 0 & 0 & 0 & 0 \\
Unidade comum & 3 & 5 & 3 & 1 & 5 & 5 & 3 & 5 & 5 & 5 & 5 & 3 & 1 \\
Processo e organização & 5 & 1 & 3 & 3 & 1 & 1 & 3 & 3 & 1 & 1 & 1 & 3 & 1 \\
Testes em separado & 5 & 3 & 3 & 5 & 1 & 3 & 5 & 3 & 3 & 3 & 3 & 3 & 5 \\
Compra de produto pronto & 3 & 3 & 1 & 1 & 0 & 0 & 1 & 0 & 3 & 0 & 0 & 1 & 5 \\
Manutenção e mantenabilidade & 5 & 5 & 3 & 3 & 1 & 1 & 5 & 1 & 3 & 1 & 1 & 3 & 3 \\
Atualização & 3 & 1 & 3 & 3 & 0 & 1 & 1 & 1 & 1 & 0 & 1 & 1 & 5 \\
Reciclagem & 5 & 3 & 3 & 3 & 1 & 1 & 5 & 3 & 3 & 0 & 1 & 3 & 3 \\
\midrule\textbf{Soma bruta}&36 & 27 & 28 & 26 & 13 & 16 & 25 & 20 & 23 & 13 & 16 & 23 & 31\\
\textbf{Posição na base}&1 & 3 & 2 & 4 & 11 & 9 & 5 & 8 & 6 & 11 & 9 & 6 & --\\
\bottomrule\end{tabularx}\endgroup
\fonte{Julgamentos anteriores à mudança para M1, preservados para rastreio. S13* não participa da base. As notas ainda não foram revalidadas para proteção direta.}
\end{table}

Na avaliação histórica, amortecimento (S01, 36) reuniu motivos fortes em fabricação, teste, manutenção e separação material. Ajuste (S03, 28) e imobilização (S02, 27) concentraram variação dimensional e reposição de desgaste. Inspeção (S04, 26) e retenção (S07, 25) exigiram combinar acesso e fechamento. Esses resultados orientam a reavaliação dos perfis sob M1; não são uma seleção das quatro maiores colunas nem importância funcional do cliente.

A planilha calcula somas e posições com empates, inclui OP-01 apenas pelo controle explícito e impede uma soma válida com nota ausente ou fora de 0/1/3/5. O registro conserva a nota histórica e sinaliza a reavaliação da fronteira M1; se a nota mudar, sinaliza também a revisão do julgamento. As somas continuam calculáveis como histórico, sem autorizar a transferência para a nova arquitetura. Não há agrupamento automático por pontuação.

Uma análise de sensibilidade substituiu 5 por 9 apenas para comparar a convenção da base com a escala consultada. As posições mudam: por exemplo, retenção passa de 25 para 37, e imobilização de 27 para 35. Os motivos para separar amortecimento, ajuste/imobilização e barreira/inspeção permanecem identificáveis; o corpo comum continua sustentado por interfaces compartilhadas. Esse teste não prova robustez a qualquer julgamento. A ordenação não deve ser usada como corte arbitrário do número de módulos.

\subsection{Arquitetura do Produto}
\label{sec:arquitetura-mim}

A análise anterior de perfis e vínculos físicos indicou quatro núcleos para C1: MD1 estrutura; MD2 ajuste e imobilização; MD3 amortecimento; MD4 contenção e inspeção. Eles permanecem hipóteses a reavaliar sob M1. MD6 é opção de evento, fora da referência. Uso e identidade do antigo MD5 foram integrados a MD1, pois cava e gravação pertencem ao corpo. A numeração de MD6 foi preservada para rastrear a opção anterior.

\begin{table}[H]
\centering\scriptsize\caption{Agrupamentos herdados a reavaliar para a proteção direta}\label{tab:mim-modulos}
\begin{tabularx}{\textwidth}{L{0.8cm}L{3cm}L{3.1cm}Y}
\toprule\textbf{ID}&\textbf{Núcleo}&\textbf{Portadores}&\textbf{Motivo e limite}\\\midrule
MD1 & Estrutura, pega e circulação & S05, S06, S08, S09, S10, S11, S12 & Predominam unidade comum e carry-over. Pega, carga, apoio, recolhimento e identidade compartilham o corpo; separar cada função aumentaria interfaces estruturais. Fechos podem ser peças de reposição internas, sem se tornarem automaticamente módulos. \\
MD2 & Ajuste e imobilização & S02, S03 & Cintas e guias atuam sobre a peça descoberta; coordenar curso, tensão, área de apoio e pressão na cerâmica. Reavaliar os motivos de agrupamento herdados. \\
MD3 & Amortecimento & S01 & Processo, teste em separado, manutenção e reciclagem são fortes. O calço tem material e vida próprios, sem exigir descarte do corpo. \\
MD4 & Contenção e inspeção & S04, S07 & Teste, manutenção, processo e fim de vida indicam separação da estrutura. O agrupamento entre visor e barreira é condicional ao fechamento contínuo e ao acesso real à cerâmica; seus perfis não são idênticos. \\
MD6* & Registro de evento opcional & S13 & Tecnologia, fornecedor, teste separado e atualização favorecem separação se OP-01 entrar no escopo. Alto escore não autoriza sua inclusão. \\
\bottomrule\end{tabularx}
\fonte{Análise própria dos perfis da MIM e dos portadores da C1; MD6* não incluído.}
\end{table}

As funções de pega, compressão, apoio, recolhimento e identidade permanecem ligadas ao corpo MD1. Fechos são partes locais de interface, sem serem contados automaticamente como módulos. S12 tem motivo de manutenção, mas separar o recolhimento como núcleo autônomo pode criar juntas no caminho de carga; a decisão permanece condicionada à estabilidade e à compressão. MD4 reúne perfis diferentes por exigir continuidade da barreira ao redor da região de inspeção. Essa integração só permanece se ambas as metas forem verificadas no mesmo conjunto.

Não se somaram escores por módulo para comparar grupos de tamanhos diferentes. A MIM não prova quantidade de módulos substituíveis. Para RP20, a hipótese é demonstrar a troca independente de MD2, MD3 e MD4. MD1 preserva a referência estrutural e a identidade da unidade; não entra nessa contagem. MD6 opcional também não foi usado para atingir três módulos.

\begin{table}[H]
\centering\scriptsize\caption{Interfaces e verificação da troca independente}\label{tab:mim-interfaces}
\begin{tabularx}{\textwidth}{L{1.1cm}L{3.7cm}YY}
\toprule\textbf{ID}&\textbf{Interface a definir}&\textbf{Troca a executar}&\textbf{Verificação após reposição}\\\midrule
I12&Base--ajuste: fixação, datum, curso e trava&Retirar/repor MD2 sem desmontar MD3/MD4&Amplitude sem troca no ajuste e deslocamento residual\\
I13&Base/ajuste--calço: assento, contato e retenção&Retirar/repor MD3 mantendo os outros núcleos&Resposta do calço e aceleração, dispersão e vida do conjunto\\
I14&Base/ajuste--barreira: fixação, vedação e acesso óptico&Retirar/repor MD4 preservando base, guias e calços&Área/defeitos verificáveis, retenção e retirada segura\\
I16*&Base--OP-01: montagem, posição, energia/leitura&Somente após inclusão no escopo&Calibração, benefício, vida e custo; sem diagnóstico de trinca\\
\bottomrule\end{tabularx}
\fonte{Plano de verificação proposto; resultados ainda não disponíveis. I16* está fora da C1 base.}
\end{table}

Executar trocas em conjunto descarregado, registrando acesso, ações, tempo e partes retiradas. Ajustar a faixa de tamanho permanece na mesma configuração: não se confunde com remover módulo por manutenção. Após cada troca, repetir as verificações afetadas e preservar a identidade. C2 e C3 exigem reavaliação da MIM para seus próprios portadores, particularmente rolamento, sensores e identificação; esta arquitetura não foi transportada automaticamente para essas concepções.
